\documentclass[11pt,a4paper,oneside]{report}

% ---------------------------------------------------------------------------
% Packages
% ---------------------------------------------------------------------------
\usepackage[T1]{fontenc}
\usepackage[utf8]{inputenc}
\usepackage{lmodern}
\usepackage[margin=2.5cm]{geometry}
\usepackage{microtype}
\usepackage[dvipsnames]{xcolor}
\usepackage{booktabs}
\usepackage{tabularx}
\usepackage{array}
\usepackage{caption}
\usepackage{tikz}
\usepackage{titlesec}
\usepackage{fancyhdr}
\usepackage{cite}
\usepackage{xurl}
\usepackage{hyperref}
\hypersetup{
  colorlinks=true,
  linkcolor=MidnightBlue,
  citecolor=MidnightBlue,
  urlcolor=MidnightBlue,
  pdftitle={Background and Related Work},
}

% ---------------------------------------------------------------------------
% Layout
% ---------------------------------------------------------------------------
\linespread{1.08}
\setlength{\parskip}{0.3em}

\definecolor{accent}{RGB}{25,55,110}

\titleformat{\chapter}[display]
  {\normalfont\bfseries\color{accent}}
  {\fontsize{48}{48}\selectfont\thechapter}
  {6pt}
  {\Huge}
  [\vspace{4pt}{\color{accent}\titlerule[0.8pt]}]
\titlespacing*{\chapter}{0pt}{0pt}{28pt}

\titleformat{\section}
  {\normalfont\Large\bfseries\color{accent}}{\thesection}{0.8em}{}
\titleformat{\subsection}
  {\normalfont\large\bfseries\color{accent}}{\thesubsection}{0.8em}{}
\titlespacing*{\section}{0pt}{1.6em}{0.7em}
\titlespacing*{\subsection}{0pt}{1.3em}{0.5em}

\pagestyle{fancy}
\fancyhf{}
\fancyhead[L]{\small\itshape Chapter~\thechapter.\ Background \& related work}
\fancyhead[R]{\small\thepage}
\renewcommand{\headrulewidth}{0.4pt}
\fancypagestyle{plain}{\fancyhf{}\fancyfoot[C]{\small\thepage}\renewcommand{\headrulewidth}{0pt}}

\captionsetup{font=small, labelfont={bf,color=accent}, labelsep=period, margin=1cm}
\captionsetup[table]{position=top, skip=6pt}

\renewcommand{\arraystretch}{1.25}
\newcolumntype{L}{>{\raggedright\arraybackslash}X}
\newcolumntype{R}{>{\raggedright\arraybackslash\bfseries}p{3.3cm}}

% Inline formatting for addresses, prefixes and example strings
\newcommand{\pfx}[1]{\texttt{#1}}

% Placeholder for figures still to be drawn
\newcommand{\figplaceholder}[1]{%
  \begin{tikzpicture}
    \draw[accent!60, dashed, line width=0.8pt, rounded corners=4pt]
      (0,0) rectangle (0.8\textwidth,5cm);
    \node[align=center, text width=0.7\textwidth, text=accent!80]
      at (0.4\textwidth,2.5cm) {\itshape Figure to be added:\\[2pt] #1};
  \end{tikzpicture}%
}

\begin{document}

\setcounter{chapter}{1}
\chapter{Background \& related work}
\label{ch:background}

% ===========================================================================
\section{Longest Prefix Match in software}
\label{sec:lpm}

Routers decide where to forward each packet using longest prefix match
(LPM)~\cite{ruiz2001survey}. Stripped of its networking context, the problem is
simple: given a set of strings $S$ and a query string $x$, find the longest
string in $S$ that is a prefix of $x$. For example, let
$S = \{\pfx{ABA}, \pfx{ABAC}\}$. For $x = \pfx{ABAS}$, the answer is \pfx{ABA},
since \pfx{ABA} is a prefix of \pfx{ABAS} but \pfx{ABAC} is not. For
$x = \pfx{ABACE}$, both are prefixes, and the answer is \pfx{ABAC} because it is
longer.

In a router, the strings are bit strings. A routing table is a set of routes,
each pairing an address prefix with a next hop. A prefix is written as an
address followed by its length in bits: \pfx{10.0.0.0/8} denotes the first
8~bits of \pfx{10.0.0.0}, i.e.\ \pfx{00001010}. For each packet, the router
finds the longest prefix in the table that matches the packet's destination
address and forwards the packet to that route's next hop.

\begin{figure}[htbp]
  \centering
  \figplaceholder{example routing table and the corresponding binary trie}
  \caption{An example routing table and the corresponding binary trie.}
  \label{fig:trie}
\end{figure}

A binary trie for IPv4 addresses needs up to 32 memory accesses per lookup,
which is usually too slow for packet forwarding. Real implementations reduce
this by matching several bits per step (multibit tries~\cite{srinivasan1999}),
by skipping chains of single-child nodes (path compression~\cite{morrison1968}),
or by indexing directly into a large table using the first bits of the
address~\cite{gupta1998}. These bring a lookup down to a handful of memory
accesses. Multibit tries and direct indexing use more memory than a binary
trie, trading it for faster lookups. With millions of packets per second to
process, packet forwarding prioritises lookup time, but these structures also
make updates more expensive, since a single route change can rewrite many
entries~\cite{ruiz2001survey}.

In software, lookup time is dominated not by computation but by memory
accesses~\cite{drepper2007}. Each step must fetch the next node, and nodes are
often not close to each other in memory, which hurts cache locality. This
matters because a cache miss can take tens to hundreds of times longer than a
hit~\cite{drepper2007}. Lookup performance depends heavily on how much of the
structure fits in cache. Cloudflare observed lookups in the Linux kernel's BPF
LPM trie slow sharply as tables grew beyond tens of thousands of
entries~\cite{fleming2025}, and Arashloo et al.\ found that a lookup table too
large for the L3 cache needed almost twice as many cores to sustain the same
packet rate~\cite{arashloo2018}. In multi-tenant settings, a software router
must keep a separate structure for every tenant, so their combined size quickly
outgrows the cache.

% ===========================================================================
\section{Multi-tenancy}
\label{sec:multitenancy}

Multi-tenancy describes a single forwarding device shared by many independent
users, called tenants. Tenants may be virtual networks in a data centre or
cloud, or customers of a virtual private network (VPN) service.

Each tenant expects its network to behave as if it were the only one, even
though it shares the underlying infrastructure with many others. Two
requirements follow. First, tenants commonly use the same private address
ranges~\cite{rfc1918}, so the same address can mean different things to
different tenants~\cite{rfc4364}. Second, tenants must be isolated, so one
tenant's routes must not affect how another tenant's packets are forwarded. For
example, two tenants may both define a route for \pfx{10.0.0.0/24}, where the
first tenant's route leads to a server in its virtual network and the second
tenant's leads to its office in London. If the second tenant adds a route
\pfx{10.0.0.5/32} for a printer in that office, packets of the first tenant
must be unaffected.

This is usually achieved by keeping a separate routing table, called a
virtual routing and forwarding (VRF) instance, for each
tenant~\cite{rfc4364}. The router identifies each packet's tenant, for example
from the interface it arrived on or from a tag carried in the packet, and
performs LPM in that tenant's VRF.

\begin{figure}[htbp]
  \centering
  \figplaceholder{one device with several VRFs; packets tagged with their tenant}
  \caption{A single forwarding device holding several VRFs. Each packet is
    tagged with its tenant and looked up in that tenant's VRF.}
  \label{fig:vrfs}
\end{figure}

Because tenants draw their addresses from the same private ranges, their tables
often contain the same prefixes, and storing them separately duplicates this
data. This adds to the memory and cache pressure described above. In the
routing tables of one cloud provider, only about 10\% of prefixes were distinct
across tenants~\cite{xu2024map}.

% ===========================================================================
\section{Multi-tenant settings}
\label{sec:settings}

We focus on three settings in which multi-tenant lookup arises: software cloud
gateways, provider-edge routers for L3VPNs, and host virtual switches. All three
forward traffic for many tenants from a single device, but they cover a wide
range of workloads. Cloud gateways serve up to a million tenants, most with very
small VRFs. Provider-edge routers serve tens to hundreds of VPN customers, each
with VRFs of thousands of routes. Host virtual switches serve tens to about a
hundred tenants on each server of a data centre. For each setting, we describe
the system, why it runs or could run in software, and its workload: the number
of tenants, the size and similarity of their tables, and how often routes change
and are looked up.

% ---------------------------------------------------------------------------
\subsection{Software cloud gateways}
\label{sec:gateways}

Large public clouds let tenants rent virtual private clouds (VPCs), which are
isolated virtual networks of virtual machines. Traffic that leaves a VPC to the
internet, to cloud services such as storage, to the tenant's own offices, or to
another VPC, passes through a cloud gateway~\cite{pan2021sailfish}. The gateway
therefore holds the complete routing state of every
tenant~\cite{pan2021sailfish}. Such gateways are commonly built in software on
ordinary x86 servers. Alibaba's gateway is built on DPDK, a framework that lets
packet-processing software bypass the operating system kernel and run directly
on the server's cores~\cite{pan2021sailfish}, \cite{dpdk}, and Google's runs
entirely in software~\cite{dalton2018andromeda}. Some providers offload part of
this traffic to programmable switches~\cite{pan2021sailfish}, but we focus on
gateways implemented in software.

Within the cloud, tenant traffic is carried using Virtual Extensible LAN
(VXLAN), which tags each packet with a VXLAN network identifier
(VNI)~\cite{rfc7348}. The gateway reads the VNI to determine which tenant the
packet belongs to, then performs LPM on the destination address in that
tenant's VRF~\cite{pan2021sailfish}.

A large cloud region has up to a million tenants~\cite{pan2021sailfish}, with
between 177K and 1.16M routing entries in total across the regions measured by
QuarkTable~\cite{lu2024quarktable}. Most tenants' tables are very small: 80\%
have fewer than 10 routes, fewer than 5\% have more than 100, and the largest
has about 1,500~\cite{lu2024quarktable}. The tables are also highly similar.
Almost 80\% of tenants draw their addresses from \pfx{10.0.0.0/8} or
\pfx{172.16.0.0/12}~\cite{lu2024quarktable}, and only about 10\% of prefixes are
distinct across tenants~\cite{xu2024map}. Prefix lengths range from \pfx{/16} to
\pfx{/32}, peaking at \pfx{/16}, \pfx{/20} and \pfx{/24}~\cite{lu2024quarktable}.

On the read side, a single x86 core processes about 1~million packets per
second, and a whole server tens of millions~\cite{pan2021sailfish}. Traffic is
heavily skewed: about 5\% of entries carry 95\% of
traffic~\cite{pan2021sailfish}, so a structure that keeps the busiest entries in
cache could benefit considerably. The tables are large relative to the cache.
Server processors provide only a few megabytes of last-level cache per core; the
Xeon used by Arashloo et al., for example, has 45\,MB shared by
18~cores~\cite{arashloo2018}. A single table of two million routes in VPP
(Vector Packet Processing), an open-source software router, already takes about
47\,MB~\cite{vppscale}, and keeping a separate table for each of a million
tenants would need far more. Updates arrive at about 1,200 per second on average
across a region (over 100~million per day)~\cite{wei2023achelous}, but often in
bursts, for example when a large tenant joins or adds many routes at
once~\cite{pan2021sailfish}. The resulting workload combines a
very large number of tenants with small, highly similar tables, traffic
concentrated on a few entries, and updates that are modest on average but
bursty.

% ---------------------------------------------------------------------------
\subsection{L3VPN provider-edge routers}
\label{sec:pe}

Here, tenants are customers of a VPN provider who want to connect their sites,
such as offices or data centres, to each other. Each site has a customer edge
(CE) router, which connects to a provider edge (PE) router maintained by the VPN
provider. The PE forwards the customer's traffic across the provider's network
using MPLS to the customer's other sites~\cite{rfc4364}. Each PE serves several
customers and keeps one VRF per customer. That VRF holds all of the customer's
VPN routes, so that any site can reach any other site~\cite{kim2008relaying}.
The size of a VRF grows with the number of sites the customer has, not with the
traffic at this PE. The PE learns each customer's routes from the customer's CE
router, either through static configuration or a routing protocol such as
BGP~\cite{rfc4364}, and exchanges them with the other PEs using multiprotocol
BGP~\cite{rfc4364}.

PEs have traditionally been hardware routers, whose forwarding tables sit in
limited on-chip memory. Kim et al.\ describe this memory as ``a key bottleneck
in provisioning new customers''~\cite{kim2008relaying}. PEs can also be
implemented in software, such as Juniper's vMX, which runs on x86 servers using
DPDK~\cite{dinham2021vmx}. In this report we focus on PEs implemented in
software.

A PE may serve from tens to hundreds of tenants; vMX, for example, supports up
to 250 VPN instances on its highest licence tier, and up to 2~million routes
(8~million at 10\,Gbps and above) on its Advance tier~\cite{dinham2021vmx}. Each
tenant may have thousands of routes. In AT\&T's network, the largest VPN
installs more than 20,000 routes in each of its VRFs~\cite{kim2008relaying}. For
example, a chain of 7,200 stores connected by a single VPN requires at least
14,400 routes in every VRF of that VPN~\cite{kim2008relaying}. At the extreme, a
VRF may contain the Internet routes themselves~\cite{rfc4364}, and some
providers carry the full Internet routing table inside a VRF in this
way~\cite{tinka2015}, which today holds about one million IPv4
prefixes~\cite{huston2025}. Tenants' tables are less similar than in the cloud
setting, since each VPN's routes lead to its own sites~\cite{luo2012tcam}, but
customers still reuse private address ranges, so some overlap remains. To our
knowledge, no published data describes the distribution of prefix lengths in
VPN tables.

On the read side, Arashloo et al.\ built a cloud provider's VPN edge whose
per-customer lookups run in software on a server, forwarding about 40\,Gbps on
10~cores~\cite{arashloo2018}. When their lookup table grew to 16~million
entries, larger than the processor's 45\,MB L3 cache, it needed 10~cores instead
of 6 for the same packet rate~\cite{arashloo2018}. Traffic is concentrated on a
small part of each table: 80\% of VRFs use only about 10\% of their prefixes, and
these active sets stay stable over weeks~\cite{kim2008relaying}. Updates are
driven by BGP. No public data on their frequency in VPNs exists, but Internet
routing tables, which are larger and more volatile, receive up to about 160
updates per second~\cite{li2014}. This suggests an upper bound for most VRFs and
is the actual rate for a VRF carrying the Internet table. Compared with cloud
gateways, PEs serve fewer tenants with larger and less similar tables; traffic
remains concentrated on a small part of each table, and updates are driven by
BGP.

% ---------------------------------------------------------------------------
\subsection{Host virtual switches}
\label{sec:vswitch}

A host virtual switch (vSwitch) is a software component running on every server
in a cloud data centre~\cite{firestone2017vfp}, \cite{dalton2018andromeda}. It
typically acts as the tunnel endpoint for the server's virtual machines,
encapsulating their packets, for example with VXLAN, before they enter the
physical network and decapsulating packets on arrival. VXLAN was originally
designed to carry layer-2 traffic between machines~\cite{rfc7348}, but in clouds
the vSwitch also routes packets at layer~3 on behalf of each
tenant~\cite{firestone2017vfp}. This leaves the physical network unaware of
tenants; multi-tenancy is handled entirely by the vSwitches.

When a vSwitch receives a packet from a tenant VM, it encapsulates it with the
VM's VNI and sends it across the physical network as an ordinary packet. To do
so, it must know the physical address of the server hosting the destination, so
it keeps an exact-match table mapping (VNI, private address) pairs to server
addresses~\cite{firestone2017vfp}, \cite{pan2021sailfish}. This table alone is
not enough, as each virtual network also has a set of routes in its VRF
describing where ranges of destination addresses should go. A virtual network
might, for example, send \pfx{10.1.0.0/16} to cloud storage, \pfx{10.2.0.0/16}
to its VPN endpoint, and \pfx{10.3.0.0/16} to other VMs in the same network. The
vSwitch therefore first performs LPM in the tenant's VRF to decide what to do
with the packet, and only then, if the destination is another VM, uses the
exact-match table to find its server~\cite{firestone2017vfp}. As a result, it
must hold a VRF for every virtual network with a VM on the server. Providers
differ in how much of this state they keep on hosts: Google installs only part
of it, sending other traffic through dedicated
gateways~\cite{dalton2018andromeda}, while Azure's vSwitch holds routing tables
directly~\cite{firestone2017vfp}. We assume the latter, so each VRF on a host
has the same contents as at the cloud gateway.

This may sound similar to the cloud gateway, but the vSwitch operates under
different constraints. Above all, it shares the server's cache and memory with
the tenants' VMs. Co-located VMs already compete for the shared last-level
cache~\cite{hu2026cacheman}, and the vSwitch adds to this contention. Memory also
has a direct cost: DRAM can account for half of a server's
cost~\cite{li2023pond}, and memory used by the vSwitch cannot be sold to
tenants. Google, for example, caps its host dataplane's memory and plans to
account for it when placing VMs~\cite{dalton2018andromeda}. In addition, VMs are
created and destroyed frequently, with most living only
minutes~\cite{hadary2020protean}, and a server may host only one or a few VMs
from each virtual network. The vSwitch must therefore add and remove whole VRFs
as VMs arrive and leave, which can be costly for structures with a large fixed
size per table.

Current server processors have up to 192 cores~\cite{amd9965}, and most cloud
VMs need only one to four cores~\cite{hadary2020protean}. A single server can
therefore host roughly 50--100 VMs, and its vSwitch may hold tens to about a
hundred VRFs, at most one per VM. Since the VRFs have the same contents as at
the gateway, most hold fewer than 10 routes~\cite{lu2024quarktable}. On the read
side, the vSwitch processes all traffic of the server's VMs, often on a small,
dedicated share of the CPU. Google reserves a single physical core for its
dataplane, which processes over 3~million packets per second within a budget of
300\,ns per packet~\cite{dalton2018andromeda}. Updates take two forms. Route
changes within a VRF arrive from the cloud's control plane, as at the gateway,
but a host sees only those for the virtual networks it serves. VRFs themselves
are added and removed as VMs arrive and leave. No published figures give the
rate per host, but with VM lifetimes of minutes~\cite{hadary2020protean}, this
happens regularly on a busy server. In summary, a vSwitch serves a modest number
of tenants with very small, likely overlapping tables, but must add and remove
VRFs dynamically and has a direct financial incentive to minimise memory use.

% ---------------------------------------------------------------------------
\subsection{Comparison}
\label{sec:comparison}

\begin{table}[htbp]
  \centering
  \small
  \caption{Workload characteristics of the three multi-tenant settings.}
  \label{tab:settings}
  \begin{tabularx}{\textwidth}{@{} R L L L @{}}
    \toprule
      & \textbf{Cloud Gateway}
      & \textbf{L3VPN PEs}
      & \textbf{Virtual host switch} \\
    \midrule
    Tenants per device
      & Up to 1M
      & Tens to hundreds
      & Tens to about a hundred \\
    Routes per tenant
      & Mostly $<$10, up to 1500
      & Thousands, more than 20,000
      & As cloud gateway \\
    Similarity between tenants
      & 10\% of prefixes are distinct
      & Lower than cloud; not measured
      & Likely high; not measured \\
    Prefix lengths
      & From 16 to 32, mostly \pfx{/16}, \pfx{/20} and \pfx{/24}
      & Customer subnets; no published distribution
      & As cloud gateway \\
    Number of active routes
      & 5\% is used by 95\% of traffic
      & 80\% of VRFs use $<$10\% of their prefixes
      & Not measured \\
    Update frequency
      & $\sim$1,200/s
      & BGP-driven, upper bound 160/s
      & N/A \\
    \bottomrule
  \end{tabularx}
\end{table}

Table~\ref{tab:settings} summarises the workloads of the three settings. Cloud
gateways combine a very large number of tenants with small, highly similar
tables and traffic concentrated on a few entries. PEs serve far fewer tenants,
but with much larger tables, and although customers still reuse private address
ranges, their tables are likely less similar, since each VPN's routes lead to
its own sites. Host virtual switches hold the same kind of small tables as cloud
gateways, but for only tens of tenants, and under tighter constraints: they
share cache and memory with the tenants' VMs, memory they use cannot be sold,
and whole VRFs must be added and removed as VMs come and go. Despite these
differences, all three share the features that motivate this project: a single
device holds many tables at once, and the tables change while traffic flows
through them. Where it has been measured, lookups also concentrate on a small
part of each table.

% ===========================================================================
\section{Concurrent updates and lock-free lookups}
\label{sec:concurrency}

To increase the throughput of forwarding in software, routers process packets on
multiple threads, each typically running on its own core~\cite{vppthreads}. This
allows the same machine to handle many packets at the same time, increasing
throughput and utilisation. Nevertheless, concurrency has its own drawbacks once
there are concurrent updates and reads of the same shared state, in our case the
forwarding table. We have to ensure that even while updates are in progress,
readers can still fetch correct information from the data structure, and that
update processing does not slow them down. This matters especially in
forwarding, where read throughput is critical and reads vastly outnumber
updates. This guarantee should hold even if threads stall for an arbitrary
time. We also have to ensure atomic visibility of updates: readers should see
the state of the data structure only before or after an update, never while it
is in progress.

In practice, updates to forwarding tables are serialised, so that at most one
writer modifies a table at a time, while many threads perform lookups
concurrently. In the Linux kernel, a single global lock (RTNL) is used for all
changes to network configuration, while readers access the shared state under
RCU without taking it~\cite{rtnetlink}. DPDK's LPM library allows lookups from
many threads in parallel, but requires that additions and deletions are not
performed by multiple threads at once~\cite{dpdkthreads}. Google's Andromeda
goes further: a control thread builds an updated copy of the flow table and
then points the packet-processing threads to it, so these threads never modify
the table and no locks are needed~\cite{dalton2018andromeda}.

How updates can affect reads is characterised by three progress guarantees, as
defined by Herlihy and Shavit~\cite{herlihy2008art}. Under a \emph{blocking}
guarantee, a stalled thread can stop the progress of others, so a read might not
be able to complete until an update has finished. Under a \emph{lock-free}
guarantee, at least one thread is always guaranteed to make progress, so a
complete stall of the whole program is impossible. Under a \emph{wait-free}
guarantee, every operation finishes in a bounded number of its own steps, so
regardless of what other threads are doing, each thread completes its
operation~\cite{herlihy1991}.

A blocking guarantee is the easiest to implement, and the usual choice for most
applications. However, blocking does not fit our requirements, as a stalled
writer could prevent all other threads from reading, reducing throughput and
slowing down readers. Preferably, readers would be wait-free, so that they are
independent of the writer, and the writer would be wait-free too, never waiting
for readers. Which guarantee readers achieve depends on the chosen memory
reclamation method, but regardless of which one is chosen, neither will block.

Giving up blocking makes memory management harder. With locks, a thread that
frees memory has exclusive access to it, so it is certain that no other thread
is reading it. Without them, even if the writer decides that part of the data
structure should be deleted, it cannot reclaim that memory immediately, as
readers might still be accessing it. Section~\ref{sec:reclamation} describes how
such memory can be reclaimed safely.

% ===========================================================================
\section{Memory reclamation methods}
\label{sec:reclamation}

Concurrent memory reclamation methods split freeing memory into two steps:
\emph{retirement} and \emph{reclamation}. First, memory is retired: it is
unlinked from the data structure, so new lookups can no longer reach it, but
lookups that started earlier may still be using it. Then it is reclaimed, once
the method can guarantee that no thread can still be accessing it, at which
point it is safe to free.

Common methods for this are reference counting, quiescent-state-based
reclamation (QSBR), epoch-based reclamation (EBR) and hazard pointers
(HP)~\cite{hart2007}. Reference counting keeps, in each node, a count of the
readers currently using it. It is simple, but inefficient, as every read must
write to shared memory. QSBR, the approach behind read-copy-update (RCU) in the
Linux kernel, waits until every reader has passed a point at which it holds no
references, guaranteeing that none can still access retired
memory~\cite{mckenney1998}. EBR maintains a global epoch, which readers announce
when they start an operation, and memory retired in an epoch is freed once no
reader can still be in that epoch~\cite{fraser2004}. With hazard pointers, each
reader publishes the pointers it is about to access in shared slots, and a node
is freed only when no thread has published a hazard pointer to
it~\cite{michael2004}.

EBR and QSBR are similar in concept: both wait for all threads to pass through
some state before deciding that memory is safe to reclaim. This is why both add
very little overhead to readers, but can leave an unbounded amount of
unreclaimed memory if a reader stalls~\cite{hart2007}. HP bounds unreclaimed
memory, but adds a cost to every read~\cite{hart2007}, \cite{michael2004}. The
two approaches also give readers different guarantees. Under EBR, a lookup only
announces the epoch and then traverses a structure of bounded depth, so it is
wait-free. Under HP, a lookup must re-check each pointer after publishing it and
retry if it has changed, so it is lock-free~\cite{michael2004}. In both cases,
the writer is wait-free, as it never waits for readers. Because EBR and QSBR are
so similar, implementing both would add little, so we consider only EBR and HP,
which sit at opposite ends of this trade-off.

% ===========================================================================
% References (typeset as a section of this chapter rather than a new chapter)
% ===========================================================================
\begingroup
\renewcommand{\bibname}{References}
\let\chapter\section
\small
\begin{thebibliography}{99}
\setlength{\itemsep}{2pt}

\bibitem{ruiz2001survey}
M.~Á. Ruiz-Sánchez, E.~W. Biersack, and W.~Dabbous, ``Survey and taxonomy of IP
address lookup algorithms,'' \emph{IEEE Network}, vol.~15, no.~2, pp.~8--23,
2001.

\bibitem{srinivasan1999}
V.~Srinivasan and G.~Varghese, ``Fast address lookups using controlled prefix
expansion,'' \emph{ACM Transactions on Computer Systems}, vol.~17, no.~1,
pp.~1--40, 1999.

\bibitem{morrison1968}
D.~R. Morrison, ``PATRICIA---Practical algorithm to retrieve information coded
in alphanumeric,'' \emph{Journal of the ACM}, vol.~15, no.~4, pp.~514--534,
1968.

\bibitem{gupta1998}
P.~Gupta, S.~Lin, and N.~McKeown, ``Routing lookups in hardware at memory access
speeds,'' in \emph{Proc. IEEE INFOCOM}, 1998, pp.~1240--1247.

\bibitem{drepper2007}
U.~Drepper, ``What every programmer should know about memory,'' Red Hat, Inc.,
Nov. 2007. [Online]. Available:
\url{https://people.redhat.com/drepper/cpumemory.pdf}

\bibitem{fleming2025}
M.~Fleming and J.~Brouer, ``A deep dive into BPF LPM trie performance and
optimization,'' \emph{The Cloudflare Blog}, Oct.~21, 2025. [Online]. Available:
\url{https://blog.cloudflare.com/a-deep-dive-into-bpf-lpm-trie-performance-and-optimization/}

\bibitem{arashloo2018}
M.~T. Arashloo, P.~Shirshov, R.~Gandhi, G.~Lu, L.~Yuan, and J.~Rexford, ``A
scalable VPN gateway for multi-tenant cloud services,'' \emph{ACM SIGCOMM
Computer Communication Review}, vol.~48, no.~1, pp.~50--55, Jan. 2018.

\bibitem{rfc1918}
Y.~Rekhter, B.~Moskowitz, D.~Karrenberg, G.~J. de~Groot, and E.~Lear, ``Address
allocation for private internets,'' RFC~1918 (BCP~5), Feb. 1996.

\bibitem{rfc4364}
E.~Rosen and Y.~Rekhter, ``BGP/MPLS IP virtual private networks (VPNs),''
RFC~4364, Feb. 2006.

\bibitem{xu2024map}
K.~Xu, D.~Jiang, Y.~Li, X.~Wang, D.~Zhang, and G.~Xie, ``MaP: Increasing node
capacity of programmable cloud gateways,'' \emph{Computer Networks}, vol.~251,
Art. no.~110590, 2024.

\bibitem{pan2021sailfish}
T.~Pan, N.~Yu, C.~Jia, J.~Pi, L.~Xu, Y.~Qiao, Z.~Li, K.~Liu, J.~Lu, J.~Lu,
E.~Song, J.~Zhang, T.~Huang, and S.~Zhu, ``Sailfish: Accelerating cloud-scale
multi-tenant multi-service gateways with programmable switches,'' in
\emph{Proc. ACM SIGCOMM}, 2021.

\bibitem{dpdk}
DPDK Project, ``Data Plane Development Kit.'' [Online]. Available:
\url{https://www.dpdk.org}

\bibitem{dalton2018andromeda}
M.~Dalton \emph{et al.}, ``Andromeda: Performance, isolation, and velocity at
scale in cloud network virtualization,'' in \emph{Proc. 15th USENIX Symposium on
Networked Systems Design and Implementation (NSDI)}, 2018, pp.~373--387.

\bibitem{rfc7348}
M.~Mahalingam, D.~Dutt, K.~Duda, P.~Agarwal, L.~Kreeger, T.~Sridhar,
M.~Bursell, and C.~Wright, ``Virtual eXtensible Local Area Network (VXLAN): A
framework for overlaying virtualized layer~2 networks over layer~3 networks,''
RFC~7348, Aug. 2014.

\bibitem{lu2024quarktable}
J.~Lu, H.~Zhao, Y.~Qi, Y.~Feng, X.~Guo, S.~Li, E.~Song, X.~Wei, B.~Lyu, R.~Wen,
and S.~Zhu, ``QuarkTable: Building compact forwarding tables for programmable
switches on public clouds,'' in \emph{Proc. 8th Asia-Pacific Workshop on
Networking (APNet)}, 2024.

\bibitem{vppscale}
FD.io, ``Scale,'' \emph{The Vector Packet Processor (VPP) v24.02
documentation}. [Online]. Available:
\url{https://s3-docs.fd.io/vpp/24.02/developer/corefeatures/fib/scale.html}

\bibitem{wei2023achelous}
C.~Wei, X.~Li, Y.~Yang, X.~Jiang, T.~Xu, B.~Yang, T.~Wu, C.~Xu, Y.~Lv, H.~Gao,
Z.~Zhang, Z.~Chen, Z.~Wang, Z.~Zhang, S.~Zhu, and W.~Chen, ``Achelous: Enabling
programmability, elasticity, and reliability in hyperscale cloud networks,'' in
\emph{Proc. ACM SIGCOMM}, 2023, pp.~769--782.

\bibitem{kim2008relaying}
C.~Kim, A.~Gerber, C.~Lund, D.~Pei, and S.~Sen, ``Scalable VPN routing via
relaying,'' in \emph{Proc. ACM SIGMETRICS}, 2008, pp.~61--72.

\bibitem{dinham2021vmx}
M.~Dinham, M.~Katti, and V.~Reddy, \emph{Day One: vMX Up and Running}, 2nd~ed.
Juniper Networks Books, 2021.

\bibitem{tinka2015}
M.~Tinka, ``[c-nsp] Internet in VRF,'' cisco-nsp mailing list, May~1, 2015.
[Online]. Available:
\url{https://puck.nether.net/pipermail/cisco-nsp/2015-May/099383.html}

\bibitem{huston2025}
G.~Huston, ``A day in the life of BGP,'' \emph{The ISP Column}, Jun. 2025.
[Online]. Available: \url{https://ispcol.potaroo.net/2025-06/daybgp.html}

\bibitem{luo2012tcam}
L.~Luo, G.~Xie, S.~Uhlig, L.~Mathy, K.~Salamatian, and Y.~Xie, ``Towards
TCAM-based scalable virtual routers,'' in \emph{Proc. 8th ACM International
Conference on Emerging Networking Experiments and Technologies (CoNEXT)}, 2012,
pp.~73--84.

\bibitem{li2014}
Y.~Li, D.~Zhang, K.~Huang, D.~He, and W.~Long, ``A memory-efficient parallel
routing lookup model with fast updates,'' \emph{Computer Communications},
vol.~38, pp.~60--71, 2014.

\bibitem{firestone2017vfp}
D.~Firestone, ``VFP: A virtual switch platform for host SDN in the public
cloud,'' in \emph{Proc. 14th USENIX Symposium on Networked Systems Design and
Implementation (NSDI)}, 2017, pp.~315--328.

\bibitem{hu2026cacheman}
X.~Hu, Y.~Cao, N.~Guan, Y.~Wu, X.~Qiu, S.~Dai, B.~Luo, S.~Cheng, F.~Qiu,
Y.~Shen, and J.~Wu, ``Cacheman: A comprehensive last-level cache management
system for multi-tenant clouds,'' in \emph{Proc. ACM SIGPLAN Symposium on
Principles and Practice of Parallel Programming (PPoPP)}, 2026.

\bibitem{li2023pond}
H.~Li, D.~S. Berger, S.~Novakovic, L.~Hsu, D.~Ernst, P.~Zardoshti, M.~Shah,
S.~Rajadnya, S.~Lee, I.~Agarwal, M.~D. Hill, M.~Fontoura, and R.~Bianchini,
``Pond: CXL-based memory pooling systems for cloud platforms,'' in \emph{Proc.
ACM International Conference on Architectural Support for Programming Languages
and Operating Systems (ASPLOS)}, 2023.

\bibitem{hadary2020protean}
O.~Hadary, L.~Marshall, I.~Menache, A.~Pan, E.~E. Greeff, D.~Dion,
S.~Dorminey, S.~Joshi, Y.~Chen, M.~Russinovich, and T.~Moscibroda, ``Protean:
VM allocation service at scale,'' in \emph{Proc. 14th USENIX Symposium on
Operating Systems Design and Implementation (OSDI)}, 2020, pp.~845--861.

\bibitem{amd9965}
AMD, ``AMD EPYC 9965,'' product specifications, 2024. [Online]. Available:
\url{https://www.amd.com/en/products/processors/server/epyc/9005-series/amd-epyc-9965.html}

\bibitem{vppthreads}
FD.io, ``Multi-threading in VPP,'' \emph{The Vector Packet Processor (VPP)
v24.06 documentation}. [Online]. Available:
\url{https://s3-docs.fd.io/vpp/24.06/developer/corearchitecture/multi_thread.html}

\bibitem{rtnetlink}
The Linux Kernel, \texttt{include/linux/rtnetlink.h} (source file; comment:
``RTNL is used as a global lock for all changes to network configuration'').
[Online]. Available:
\url{https://android.googlesource.com/kernel/common/+/cb0a3edf8d00/include/linux/rtnetlink.h}

\bibitem{dpdkthreads}
DPDK Project, ``Thread safety of DPDK functions,'' in \emph{DPDK Programmer's
Guide}, v17.11. [Online]. Available:
\url{https://dpdk.readthedocs.io/en/v17.11/prog_guide/thread_safety_dpdk_functions.html}

\bibitem{herlihy2008art}
M.~Herlihy and N.~Shavit, \emph{The Art of Multiprocessor Programming}. Morgan
Kaufmann, 2008.

\bibitem{herlihy1991}
M.~Herlihy, ``Wait-free synchronization,'' \emph{ACM Transactions on
Programming Languages and Systems}, vol.~13, no.~1, pp.~124--149, 1991.

\bibitem{hart2007}
T.~E. Hart, P.~E. McKenney, A.~D. Brown, and J.~Walpole, ``Performance of memory
reclamation for lockless synchronization,'' \emph{Journal of Parallel and
Distributed Computing}, vol.~67, no.~12, pp.~1270--1285, 2007.

\bibitem{mckenney1998}
P.~E. McKenney and J.~D. Slingwine, ``Read-copy update: Using execution history
to solve concurrency problems,'' in \emph{Proc. Parallel and Distributed
Computing and Systems (PDCS)}, Las Vegas, NV, USA, Oct. 1998, pp.~509--518.

\bibitem{fraser2004}
K.~Fraser, ``Practical lock-freedom,'' Ph.D. dissertation, University of
Cambridge, 2004. Published as Technical Report UCAM-CL-TR-579.

\bibitem{michael2004}
M.~M. Michael, ``Hazard pointers: Safe memory reclamation for lock-free
objects,'' \emph{IEEE Transactions on Parallel and Distributed Systems},
vol.~15, no.~6, pp.~491--504, 2004.

\end{thebibliography}
\endgroup

\end{document}
